\subsection*{\theappendixitem.\ 局部环的归纳极限}
\phantomsection
\addcontentsline{toc}{subsection}{\theappendixitem.\ 局部环的归纳极限}

设 I 是一个非空的右过滤预序集，\((\mathbf{A}_{\alpha}, \varphi_{\beta\alpha})\) 是关于 I 的一个环的归纳系统。假设对每个 \(\alpha \in I\) ，环 \(A_{\alpha}\) 都是局部的，其极大理想为 \(m_{\alpha}\) ，诸同态 \(\varphi_{\beta\alpha}\) 都是局部且平坦的，并且有 \(\varphi_{\beta\alpha}(m_{\alpha}) A_{\beta} = m_{\beta}\) 对 \(\beta \geqslant \alpha\) 成立。用 A 表示诸 \(A_{\alpha}\) 的归纳极限，且对每个 \(\alpha \in I\) ，设 \(\varphi_{\alpha}: A_{\alpha} \to A\) 为典范同态。

命题 1. --- a) 环 A 是局部的，其极大理想为 \(\mathfrak{m} = \varinjlim \mathfrak{m}_{\alpha}\) 。对每个 \(\alpha \in \mathbf{I}\) ，同态 \(\varphi_{\alpha}\) 是局部且平坦的，且有 \(\varphi_{\alpha}(\mathfrak{m}_{\alpha})\) A = \(\mathfrak{m}\) 。

\begin{enumerate}
\def\labelenumi{\alph{enumi})}
\setcounter{enumi}{1}
\item
  若 \(\mathbf{A}_{\alpha}\) 对每个 \(\alpha \in \mathbf{I}\) 都是 Noether 的，则 \(\mathbf{A}\) 是 Noether 的。
\item
  令 \(m = \varinjlim m_{\alpha}\) ；它是 A 的一个理想。商环 A/m 是诸域 \(A_{\alpha}/m_{\alpha}\) 的归纳极限，因此是一个域（A，I，第 116 页，命题 3）。此外，A - m 的每个元素在 A 中可逆：事实上，设 \(x \in A - m\) ；存在 \(\alpha \in I\) 和 \(\xi \in A_{\alpha}\) 使得 \(x = \varphi_{\alpha}(\xi)\) ；我们有 \(\xi \notin m_{\alpha}\) ，因而 \(\xi\) 在 \(A_{\alpha}\) 中可逆，而 x 在 A 中可逆。于是，A 是一个以 m 为极大理想的局部环。设 \(\alpha \in I\) 。由关系式 \(\varphi_{\beta\alpha}(m_{\alpha}) A_{\beta} = m_{\beta}\) （\(\beta \geqslant \alpha\) ）取归纳极限即得 \(\varphi_{\alpha}(m_{\alpha}) A = m\) ；最后，由 I，§ 2，第 7 目，命题 9，同态 \(\varphi_{\alpha}\) 是平坦的。
\item
  设 \(\hat{\mathbf{A}}\) 是 \(\mathbf{A}\) 关于 m-进拓扑的分离完备化，\(\pi\) 是从 \(\mathbf{A}\) 到 \(\hat{\mathbf{A}}\) 的典范映射。假设诸环 \(\mathbf{A}_{\alpha}\) 都是 Noether 的。固定 \(\alpha \in \mathbf{I}\) ，证明环 \(\hat{\mathbf{A}}\) 在 \(\mathbf{A}_{\alpha}\) 上是 Noether 且平坦的。由假设，\(m_{\alpha}\) 是 \(\mathbf{A}_{\alpha}\) 的一个有限生成理想，故 \(m = \varphi_{\alpha}(m_{\alpha})\) 。\(m\) 是 \(\mathbf{A}\) 的一个有限生成理想。由此得出 \(\hat{m}\) （\(\hat{\mathbf{A}}\) 的极大理想）等于 \(m\hat{\mathbf{A}}\) （III，§ 2，第 12 目，命题 16 的推论 2 以及第 13 目，命题 19），从而是有限生成的。于是，环 \(\hat{\mathbf{A}}\) 是 Noether 的（同前，第 10 目，定理 2 的推论 5）。另一方面，对每个 \(n \in \mathbf{N}\) ，商 \(\hat{\mathbf{A}}/\hat{m}^{n}\) 同构于 \(\mathbf{A}/m^{n}\) （同前，第 12 目，命题 16 的推论 2 以及公式 (21)），这就意味着 \(\hat{\mathbf{A}}/\pi \circ \varphi_{\alpha}(m_{\alpha}^{n})\) \(\hat{\mathbf{A}}\) 同构于 \(\mathbf{A} \otimes_{\mathbf{A}_{\alpha}} (\mathbf{A}_{\alpha}/m_{\alpha}^{n})\) ；由于 \(\mathbf{A}\) 是平坦 \(\mathbf{A}_{\alpha}\)-模，\((\mathbf{A}_{\alpha}/m_{\alpha}^{n})\) -模 \(\hat{\mathbf{A}}/\pi \circ \varphi_{\alpha}(m_{\alpha}^{n})\) \(\hat{\mathbf{A}}\) 对每个 \(n \in \mathbf{N}\) 都是平坦的。由 III，§ 5，第 4 目，命题 2，\(\mathbf{A}_{\alpha}\) -模 \(\hat{\mathbf{A}}\) 关于 \(m_{\alpha}\) 是理想分离的；由同前，第 2 目，定理 1，\(\mathbf{A}_{\alpha}\) -模 \(\hat{\mathbf{A}}\) 因而是平坦的。取归纳极限即得 \(\hat{\mathbf{A}}\) 在 \(\mathbf{A}\) 上是（忠实）平坦的（I，§ 2，第 7 目，命题 9），从而 \(\mathbf{A}\) 是 Noether 的（I，§ 3，第 5 目，命题 8 的推论）。
\end{enumerate}

\refstepcounter{appendixitem}
